% Part: incompleteness
% Chapter: representability-in-q
% Section: representing-relations

\documentclass[../../../include/open-logic-section]{subfiles}

\begin{document}

\olfileid{inc}{req}{rel}

\olsection{表示关系}

下面说明\emph{关系}“可表示”的含义。

\begin{defn}
\ollabel{defn:representing-relations} 自然数上的关系 $R(x_0,\dots,x_k)$
是{\em 在 $\Th{Q}$ 中可表示的}，如果存在公式 $!A_R(x_0,\dots,x_k)$，使得
每当 $R(n_0,\dots,n_k)$ 为真时，$\Th{Q}$ 证明 $!A_R(\num{n_0},\dots,\num{n_k})$，
并且每当 $R(n_0,\dots,n_k)$ 为假时，$\Th{Q}$ 证明 $\lnot !A_R(\num{n_0},
\dots, \num{n_k})$。
\end{defn}

\begin{thm}
\ollabel{thm:representing-rels} 一个关系在 $\Th{Q}$ 中可表示，
当且仅当它是可计算的。
\end{thm}

\begin{proof}
证明正向：假设 $R(x_0,\dots,x_k)$ 由公式 $!A_R(x_0,\dots,x_k)$ 所表示。
以下是计算 $R$ 的算法：
对输入 $n_0$, \dots,~$n_k$，同时搜索 $!A_R(\num{n_0}, \dots, \num{n_k})$ 的证明和
$\lnot !A_R(\num{n_0}, \dots, \num{n_k})$ 的证明。
由假设，这个搜索必定会找到其中一个证明；
如果找到的是第一个，就报告“是”，否则报告“否”。

反之，假设 $R(x_0, \dots, x_k)$ 是可计算的。
根据定义，这意味着函数 $\Char{R}(x_0, \dots, x_k)$ 是可计算的。
根据 \olref[int]{thm:representable-iff-comp}，$\Char{R}$
可由某个公式表示，记为 $!A_{\Char{R}}(x_0, \dots, x_k,
y)$。令 $!A_R(x_0, \dots, x_k)$ 为公式 $!A_{\Char{R}}(x_0,
\dots, x_k, \num{1})$。
那么对任意 $n_0$, \dots,~$n_k$，
如果 $R(n_0,
\dots, n_k)$ 为真，则 $\Char{R}(n_0, \dots, n_k) = 1$，此时 $\Th{Q}$ 证明
$!A_{\Char{R}}(\num{n_0}, \dots, \num{n_k},
\num{1})$，从而 $\Th{Q}$ 证明 $!A_R(\num{n_0}, \dots,
\num{n_k})$。
另一方面，如果 $R(n_0, \dots, n_k)$ 为假，则 $\Char{R}(n_0, \dots, n_k) = 0$。
这意味着 $\Th{Q}$ 证明 \[
\lforall[y][(!A_{\Char{R}}(\num{n_0}, \dots, \num{n_k}, y) \lif y =
  \num{0})].
\]
由于 $\Th{Q}$ 证明 $\eq/[\num{0}][\num{1}]$，$\Th{Q}$ 证明 $\lnot !A_{\Char{R}}(\num{n_0}, \dots, \num{n_k}, \num{1})$，
从而 $\Th{Q}$ 证明 $\lnot !A_R(\num{n_0}, \dots, \num{n_k})$。
\end{proof}

\begin{prob}
证明：如果 $R$ 在~$\Th{Q}$ 中可表示，那么~$\Char{R}$ 也可表示。
\end{prob}

\end{document}